\subsection{方案 1：本地工作流 (TeX Live + VS Code)}

\subsubsection{编译环境安装}
\begin{itemize}
    \item \textbf{Mac OS}: 终端执行 \texttt{brew install {-}{-}cask mactex-no-gui}
    \item \textbf{Windows}: 官网下载 \texttt{install-tl-windows.exe}，执行完整安装。
    \item \textbf{Ubuntu}: 终端执行 \texttt{sudo apt install texlive-full}
\end{itemize}

\subsubsection{编辑器配置}
推荐使用 VS Code 配合 LaTeX Workshop 插件。
\begin{itemize}
    \item \textbf{VS Code 设置 (\texttt{.vscode/settings.json})}: 
    配置 recipes 编译链为 \texttt{xelatex $\rightarrow$ bibtex $\rightarrow$ xelatex*2}。为了节省性能，建议关闭自动编译，设置 \texttt{"latex-workshop.latex.autoBuild.run": "never"}。
    \item \textbf{快捷键配置 (\texttt{keybindings.json})}: 
    建议将 \texttt{Cmd+S} (或 Ctrl+S) 绑定为 \texttt{latex-workshop.build} 命令。
\end{itemize}

\subsubsection{核心编译指令}
标准的编译指令为 \texttt{xelatex main.tex}。复杂文档包含参考文献时需配合 bibtex。

\subsection{方案 2：云端工作流 (Overleaf)}
适合不方便配置本地环境或需要多人协作的用户。
\begin{enumerate}
    \item 注册并登录 Overleaf。
    \item 点击 New Project。
    \item 编写代码后按 \texttt{Cmd+Enter}（或 Ctrl+Enter）进行编译。
    \item \textbf{中文支持}: 需在文档导言区加入 \texttt{\textbackslash usepackage[UTF8]\{ctex\}}。
\end{enumerate}
